\section{What is Motion?}

One of the most fundamental questions in physics is:

\begin{center}
	\emph{What does it mean for an object to move?}
\end{center}

Suppose we observe an electron.

At time

\[
t_1
\]

its position is

\[
\mathbf r_1
=
(2,1,0)\ {\rm nm}.
\]

A short time later,

\[
t_2,
\]

its position becomes

\[
\mathbf r_2
=
(2.3,1.1,0)\ {\rm nm}.
\]

The position has changed.

Therefore,

\begin{keyidea}
Motion is the change of position with time.
\end{keyidea}

Motion is therefore not a property of space alone.

It requires two ingredients:

\begin{itemize}
	\item position,
	\item time.
\end{itemize}

%%%%%%%%%%%%%%%%%%%%%%%%%%%%%%%%%%%%%%%%%%%%%%%%%%%%%%%%%%%%%%%%%%%%%%%%%%%%%%%%
